% TOP_PROBLEM: {"id": 30007491, "title": "Uniform quasipolynomial worst-case trace reconstruction", "category_id": 15, "category": {"id": 15, "name": "computer_science", "display_name": "Computer Science", "description": "Computational complexity, algorithms, and theoretical CS.", "slug": "computer-science", "order_index": 15, "created_at": "2026-07-31T15:26:25.671Z"}, "status": "open"}
% FIRST_POSTED: 2026-10-06
% ATTEMPT_STATUS: unresolved
% OPENAI_CLAIM: {"schema_version":1,"source_repo":"https://github.com/openai/math","source_commit":"adc7f1241b42e322a6451854ab7e4b4c146bf78a","release_date":"2026-10-06","match":"partial","resolution":"partial","independently_reviewed":false,"second_pass":"n/a","families":[{"family":"122","title":"Quantitative trace-reconstruction bounds with a uniform decoder","subject":"Theoretical computer science","manuscripts":[{"title":"Uniform quasipolynomial-time trace reconstruction","pdf_path":"preprints/Uniform-quasipolynomial-time-trace-reconstruction-October-5-2026/uniform-trace-reconstruction.pdf","readme_path":"preprints/Uniform-quasipolynomial-time-trace-reconstruction-October-5-2026/README.md","date":"2026-10-05","theorem_ref":"thm:main (manuscript/main.tex:78)","statement_sources":[{"path":"preprints/Uniform-quasipolynomial-time-trace-reconstruction-October-5-2026/build/manuscript/main.tex","line":78,"label":"thm:main"}]}],"lean":{"doc_path":"lean/docs/122.md","comparators":[],"declarations":[],"covers_main_theorem":false},"reasoning_trace":null}]}
\documentclass[11pt]{article}
\usepackage[margin=25mm]{geometry}
\usepackage{fontspec}
\setmainfont{Latin Modern Roman}
\usepackage{amsmath,amssymb,mathtools}
\usepackage{unicode-math}
\setmathfont{Latin Modern Math}
\usepackage{enumitem}
\usepackage{xurl,hyperref,bookmark}
\hypersetup{colorlinks=true,urlcolor=blue}
\setcounter{secnumdepth}{0}
\setlength{\parindent}{0pt}
\setlength{\parskip}{5pt}
\setlength{\emergencystretch}{3em}
\begin{document}
\section*{Uniform quasipolynomial worst-case trace reconstruction}

\subsection{OpenAI's claim}
OpenAI presents the following claims relevant to this question, with manuscript numbers following paper order in Sources. In manuscript 1 (family 122; 2026-10-05), the asserted result is that every binary word can be reconstructed with known rational retention probability using a uniform algorithm with quasipolynomial sample and bit bounds at fixed retention, and polynomial bounds for sufficiently small deletion.

\subsection{Mathematical statement of the claimed result}
\paragraph{Manuscript 1.}
For a binary word $x\in\{0,1\}^n$, $n\ge2$, an ordinary deletion trace independently retains each symbol with rational probability $p=a/b\in(0,1]$ and preserves order. Put $q=1-p$, let $L$ be the bit length of the reduced numerator and denominator, and define
\[
 X(n,p)=\begin{cases}1+\dfrac{\log n}{1+\log(1/q)},&q>0,\\1,&q=0,\end{cases}
 \qquad B(n,p)=p^{-1}(\log n)X(n,p)^2(1+\log X(n,p))^6.
\]
The claim gives absolute $C,c\ge1$ and one randomized multitape Turing machine, given $n,p$ in binary, that requests at most
\[
 M_C(n,p)=\lceil\exp(CB(n,p))\rceil
\]
independent traces. For each input word $x$, reconstruction succeeds with probability at least $2/3$. Regardless of the internal random tape and of the received stream, provided all its strings have length bounded by $n$, the running time is at most
\[
 (n+L+M_C(n,p))^c
\]
bit operations, including preprocessing, finite-precision calculations, received bits and output. Only external generation of independent traces is treated as an oracle operation.

\subsection{Outline of the argument}
For manuscript 1 (family 122; 2026-10-05), low-order trace moments propagate a distinguishing signal through scales using damped products and prescribed frequencies. Positive moment constraints recover prefixes, and a bounded feasibility iteration turns the statistical separation into an explicit finite-precision decoder.

\subsection{Scope relative to this problem}
This is an algorithmic reconstruction claim with known binary length and retention data, beyond a sample-only existence estimate. Complete coverage of the exact uniform bit-time and sample guarantee is not established by this catalogue, so the claim is partial.

\subsection{Formal verification}
No Lean formalization is listed at this snapshot for manuscript 1 (family 122; 2026-10-05).

This site has not executed the Lean code. Described statement coverage is not an independent proof check.

\subsection{Status caveat}
The manuscripts were produced by an internal OpenAI model. OpenAI cautions that results without formalization can contain mistakes. This site has not independently checked the proofs; the label records the scope of OpenAI's claim.

\subsection{Sources}
\begin{itemize}
\item OpenAI, \emph{Uniform quasipolynomial-time trace reconstruction}, 2026-10-05; citation key \texttt{OAI:Uniform-quasipolynomial-time-trace-reconstruction-October-5-2026}. \href{https://github.com/openai/math/blob/adc7f1241b42e322a6451854ab7e4b4c146bf78a/preprints/Uniform-quasipolynomial-time-trace-reconstruction-October-5-2026/uniform-trace-reconstruction.pdf}{PDF (GitHub)}; \href{https://github.com/openai/math/raw/adc7f1241b42e322a6451854ab7e4b4c146bf78a/preprints/Uniform-quasipolynomial-time-trace-reconstruction-October-5-2026/uniform-trace-reconstruction.pdf}{PDF (direct)}; \href{https://github.com/openai/math/blob/adc7f1241b42e322a6451854ab7e4b4c146bf78a/preprints/Uniform-quasipolynomial-time-trace-reconstruction-October-5-2026/README.md}{Manuscript page and citation}; \href{https://github.com/openai/math/blob/main/preprints/Uniform-quasipolynomial-time-trace-reconstruction-October-5-2026/uniform-trace-reconstruction.pdf}{Latest version}; Statement source: \href{https://github.com/openai/math/blob/adc7f1241b42e322a6451854ab7e4b4c146bf78a/preprints/Uniform-quasipolynomial-time-trace-reconstruction-October-5-2026/build/manuscript/main.tex\#L78}{thm:main}.
\item Family 122: \href{https://github.com/openai/math/blob/adc7f1241b42e322a6451854ab7e4b4c146bf78a/lean/docs/122.md}{Lean scope}.
\item \href{https://github.com/openai/math/blob/adc7f1241b42e322a6451854ab7e4b4c146bf78a/CONTENTS.md}{OpenAI catalogue entry}.
\end{itemize}
\end{document}
